% ============================================================
%  ch05 —— §3 GL_n 的例（原文 p.19–p.22）
% ============================================================
\routenote{p.19--p.22}

\section{例：$GL_n$ 的 Kloosterman 层}

本节对 $\PP^1$ 上的常群 $G=GL_n$ 计算 Kloosterman 层 $\Kl_{GL_n}(\varphi,\chi)$。它的 Langlands 对偶是 $GL_{n,\mathbb{Q}_\ell}$，记标准表示为 $\mathrm{Std}$。于是描述 $\PP^1\setminus\{0,\infty\}$ 上的这个 $GL_n$-局部系统，就等同于描述秩 $n$ 的局部系统 $\Kl^{\mathrm{Std}}_{GL_n}(\varphi,\chi)$。我们将看到这个秩 $n$ 局部系统与 Deligne 在 [10] 中定义的经典 Kloosterman 层一致。

\subsection{另一种模解释}

我们想用向量丛来解释 $GL_n(1,2)$-丛。先定义 $\Bun_{GL_n(1,2)}$ 的一个变体。设 $\mathrm{Bun}_{n,1,2}$ 为分类如下数据的叠 $(E,F^*E,\{v_i\},F_*E,\{v_i\})$：
\begin{enumerate}[label=(\arabic*), leftmargin=1.9em, itemsep=.15em]
\item $E$ 是 $\PP^1$ 上秩 $n$ 的向量丛；
\item 一个递减滤过 $F^*E$，给出 $E$ 在 $0$ 处纤维的完全旗：
\[
E=F^0E\supset F^1E\supset\cdots\supset F^nE=E(-\{0\});
\]
\item 对每个 $i=1,\cdots,n$，一个非零向量 $v_i\in F^{i-1}E/F^iE$；
\item 一个递增滤过 $F_*E$，给出 $E$ 在 $\infty$ 处纤维的完全旗：
\[
E(-\{\infty\})=F_0E\subset F_1E\subset\cdots\subset F_nE=E;
\]
\item 对 $i=1,\cdots,n$，一个向量 $v_i\in F_iE/F_{i-2}E$，它不落在 $F_{i-1}E/F_{i-2}E$ 中（这里约定 $F_{-1}E=(F_{n-1}E)(-\{\infty\})$）。
\end{enumerate}
注意 $\mathrm{Bun}_{n,1,2}$ 是 $\PP^1$ 上 $G$-挠子的模叠，其中 $\mathbb{G}$ 是 $\PP^1$ 上的 Bruhat--Tits 群概形，满足
\begin{itemize}[leftmargin=1.6em, itemsep=.15em]
\item $\mathbb{G}|_{\PP^1\setminus\{0,\infty\}}=GL_n\times(\PP^1\setminus\{0,\infty\})$；
\item $\mathbb{G}(\mathcal O_0)=I_{GL_n}(1)^{\mathrm{opp}}$；
\item $\mathbb{G}(\mathcal O_\infty)=Z_{GL_n}(1)(k[[s]])\cdot I_{SL_n}(2)\supset I_{GL_n}(2)$。
\end{itemize}
$\mathbb{G}$ 与 $GL_n(1,2)$ 的唯一差别在于：它们在 $\infty$ 处对中心 $\mathbb{G}_m=Z_{GL_n}$ 取不同的水平结构。因此有自然态射 $GL_n(1,2)\to\mathbb{G}$，从而有自然态射 $\Bun_{GL_n(1,2)}\to\mathrm{Bun}_{n,1,2}$，它是 $\mathbb{G}_a$-挠子。

选取 $E$ 在 $\PP^1\setminus\{0,\infty\}$ 上的平凡化，可以把 $\mathrm{Bun}_{n,1,2}$ 的模问题改写如下。设 $\Lambda$ 是以 $\{e_1,\cdots,e_n\}$ 为基的自由 $k[t,t^{-1}]$-模。令 $e_{i+jn}=t^je_i$（$1\le i\le n$，$j\in\mathbb{Z}$），则 $\Lambda$ 是以 $\{e_i\}_{i\in\mathbb{Z}}$ 为基的 $k$-向量空间。对任意 $k$-代数 $R$，$R\otimes\Lambda$ 中的 $R[t]$-格是指 $R[t]$-子模 $\Lambda'\subset R\otimes\Lambda$，使得存在 $M\in\mathbb{Z}_{>0}$ 满足
\[
\operatorname{Span}_R\{e_i\mid i>M\}\subset\Lambda'\subset\operatorname{Span}_R\{e_i\mid i\ge -M\},
\]
且 $\Lambda'/\operatorname{Span}_R\{e_i\mid i>M\}$ 与 $\operatorname{Span}_R\{e_i\mid i\ge-M\}/\Lambda'$ 都是投射 $R$-模。类似地可定义 $R\otimes\Lambda$ 中的 $R[t^{-1}]$-格。

设 $\widehat{\mathrm{Bun}}_{n,1,2}(R)$ 分类数据 $(\Lambda_*,\{\mathbf v_i\}_{0\le i\le n-1},\Lambda^*,\{\mathbf v_i\}_{1\le i\le n})$：
\begin{enumerate}[label=(\arabic*), leftmargin=1.9em, itemsep=.15em]
\item $R\otimes\Lambda\supset\Lambda_0\supset\Lambda_1\supset\cdots\supset\Lambda_n=t\Lambda_0$ 是 $R[t]$-格链，使得 $\Lambda_i/\Lambda_{i+1}$ 是秩 $1$ 的投射 $R$-模。对任意 $1\le i\le n$、$j\in\mathbb{Z}$，令 $\Lambda_{i+jn}=t^j\Lambda_i$。
\item $\mathbf v_i\in\Lambda_i/\Lambda_{i+1}$ 是 $R$-基，$i=0,\cdots,n-1$；
\item $R\otimes\Lambda\supset\Lambda^n\supset\Lambda^{n-1}\supset\cdots\supset\Lambda^0=t^{-1}\Lambda^n$ 是 $R[t^{-1}]$-格链，使得 $\Lambda^i/\Lambda^{i+1}$ 是秩 $1$ 的投射 $R$-模。对任意 $1\le i\le n$、$j\in\mathbb{Z}$，令 $\Lambda^{i+jn}=t^j\Lambda^i$。
\item $\mathbf v_i\in\Lambda^i/\Lambda^{i-2}$，它在 $\Lambda^i/\Lambda^{i-1}$ 中的像是 $R$-基，$i=1,\cdots,n$。
\end{enumerate}
群 $GL_n(k[t,t^{-1}])$（视作 $k$ 上的 ind-群概形）作用在 $\Lambda$ 上，从而作用在叠 $\widehat{\mathrm{Bun}}_{n,1,2}$ 上。此外，$\mathrm{Bun}_{n,1,2}$ 自然同构于商叠 $\widehat{\mathrm{Bun}}_{n,1,2}/GL_n(k[t,t^{-1}])$。

与格链 $(\Lambda_*,\Lambda^*)$ 相伴的是 $\operatorname{Spec}R$ 上局部常值的整值函数
\[
\deg(\Lambda_*,\Lambda^*):=\chi_R\Big(\Lambda_0\oplus\Lambda^0\xrightarrow{\ (\iota_0,\iota^0)\ }R\otimes\Lambda\Big),
\]
其中 $\iota_0,\iota^0$ 是包含映射，右端的 Euler 特征定义为 $\operatorname{rk}_R\ker(\iota_0,\iota^0)-\operatorname{rk}_R\operatorname{coker}_R(\iota_0,\iota^0)$。

对 $d\in\mathbb{Z}$，设 $\mathrm{Bun}^d_{n,1,2}\subset\mathrm{Bun}_{n,1,2}$ 为分类满足 $\deg(\Lambda_*,\Lambda^*)=d$ 的 $\{\Lambda_*,\mathbf v_i,\Lambda^*,\mathbf v_i\}$ 的子叠。

大胞腔的嵌入 $j_d:T\times\mathbb{G}_a^n\hookrightarrow\mathrm{Bun}^d_{n,1,2}$ 可如下固定。对 $a=(a_1,\cdots,a_n)\in T(R)=(R^\times)^n$ 与 $b=(b_1,\cdots,b_n)\in R^n$，点 $j_d(a,b)\in\mathrm{Bun}^d_{n,1,2}(R)$ 由如下数据的 $GL_n(R[t,t^{-1}])$-轨道给出：
\begin{enumerate}[label=(\arabic*), leftmargin=1.9em, itemsep=.15em]
\item $\Lambda_i=\operatorname{Span}_R\{e_j\mid j>i\}\subset R\otimes\Lambda$；
\item $\mathbf v_i=a_{i+1}e_{i+1}\in\Lambda_i/\Lambda_{i+1}$；
\item $\Lambda^i=\operatorname{Span}_R\{e_j\mid j\le i+d\}\subset R\otimes\Lambda$；
\item $\mathbf v_i=e_{i+d}+b_ie_{i+d-1}\in\Lambda^i/\Lambda^{i-2}$。
\end{enumerate}
特别地，有基点 $\star_d=j_d(1,0)\in\mathrm{Bun}^d_{n,1,2}$。记 $\star_d$ 的 $k[t]$-链（相应地 $k[t^{-1}]$-链）为 $\Lambda_*(\star_d)$（相应地 $\Lambda^*(\star_d)$）。注意任意点 $j_d(a,b)$ 的底向量丛是丛
\[
E_d=\mathcal O(m+1)^{\oplus r}\oplus\mathcal O(m)^{\oplus n-r},
\]
其中 $m\in\mathbb{Z}$、$0\le r\le n-1$ 由 $d=mn+r$ 唯一确定。

\subsection{与 $GL_n$ 的标准表示相伴的 Kloosterman 层}

设 $\chi:T(k)\to\mathbb{Q}_\ell^\times$ 是特征，它在 $T$ 上定义 Kummer 局部系统 $\mathrm{Kum}_\chi$。由定理 1，$\mathrm{Bun}^d_{n,1,2}$ 上的 perverse 层
\[
A_{\varphi,\chi}^d=j_{d,!}\big(\mathrm{Kum}_\chi[n]\boxtimes\varphi^*\mathrm{AS}_\psi[n]\big)
\]
构成 $\mathrm{Bun}_{n,1,2}$ 上的 Hecke 特征层。

考虑由余权 $\omega_1=(1,0,\cdots,0)$ 给出的 Hecke 对应 $\Hecke_{\omega_1}$（$\omega_1$ 定义表示 $\mathrm{Std}$）。把它限制到 $\mathrm{Bun}^0$ 与 $\mathrm{Bun}^1$ 的分支以及曲线 $\PP^1\setminus\{0,\infty\}$ 上：
\begin{equation}
\begin{tikzcd}[column sep=huge, row sep=large]
& \Hecke_{\omega_1} \arrow[dl, "\mathrm{pr}_1"'] \arrow[dr, "\mathrm{pr}_2"] & \\
\mathrm{Bun}^0_{n,1,2} & & \mathrm{Bun}^1_{n,1,2}\times(\PP^1\setminus\{0,\infty\})
\end{tikzcd}
\end{equation}
我们想计算 $\PP^1\setminus\{0,\infty\}$ 上由下式刻画的特征值局部系统 $\Kl^{\mathrm{Std}}_{GL_n}(\varphi,\chi)$：
\[
\mathrm{pr}_{2,!}\mathrm{pr}_1^*A^0_{\varphi,\chi}[n-1]\cong A^1_{\varphi,\chi}\boxtimes\Kl^{\mathrm{Std}}_{GL_n}(\varphi,\chi).
\]
$\Hecke_{\omega_1}$ 上的相交上同调层 $\mathrm{IC}_{\omega_1}$ 就是 $\mathbb{Q}_\ell[n-1]$，因为 $\mathrm{pr}_1:\Hecke_{\omega_1}\to\mathrm{Bun}_{n,1,2}$ 是 $\mathbb{P}^{n-1}$-丛\footnote{在这个计算中，我们不把 $\mathrm{IC}_{\omega_1}$ 归一化为权 $0$，如注 2.9 那样。故严格说来，层 $\Kl^{\mathrm{Std}}_{GL_n}(\varphi,\chi)$ 与 §2.3 中定义的那个相差一个 Tate 扭转 $\big(\tfrac{n-1}{2}\big)$。}。把图 (3.1) 限制到 $\star_1\times(\PP^1\setminus\{0,\infty\})\subset\mathrm{Bun}^1_{n,1,2}\times(\PP^1\setminus\{0,\infty\})$ 在 $\mathrm{pr}_2$ 下的纤维，以及大胞腔在 $\mathrm{pr}_1$ 下的纤维 $\mathrm{GR}^\circ_{\omega_1}$，得到
\begin{equation}
\begin{tikzcd}[column sep=large, row sep=large]
\mathrm{GR}^\circ_{\omega_1} \arrow[r, hook] \arrow[d, "\mathrm{pr}_1"'] \arrow[dr, "\pi"] & \mathrm{GR}_{\omega_1} \arrow[d, "\mathrm{pr}_1"] & \\
T\times\mathbb{G}_a^n \arrow[r, "j_0"'] & \mathrm{Bun}^0_{n,1,2} & \PP^1\setminus\{0,\infty\}
\end{tikzcd}
\end{equation}
开子概形 $\mathrm{GR}^\circ_{\omega_1}\subset\mathrm{GR}_{\omega_1}$ 的 $R$-点分类 $R[t,t^{-1}]$-同态 $M:R\otimes\Lambda\to R\otimes\Lambda$，它把链 $(\Lambda_*(\star_0),\Lambda^*(\star_0))$ 送到 $(\Lambda_*(\star_1),\Lambda^*(\star_1))$，模去左复合以 $(R\otimes\Lambda,\Lambda_*(\star_0),\Lambda^*(\star_0))$ 的自同构。关于 $R\otimes\Lambda$ 的 $R[t,t^{-1}]$-基 $\{e_1,\cdots,e_n\}$，任何这样的 $M$ 形如
\[
\begin{pmatrix}
a_1 & tb_n & & \\
b_1 & a_2 & & \\
& b_2 & a_3 & \\
& & \ddots & \ddots\\
& & & b_{n-1} & a_n
\end{pmatrix},
\]
模去右乘对角矩阵（这里 $a_i,b_i\in R^\times$）。因此可以把 $M$ 的矩阵归一化为
\begin{equation}
M=\begin{pmatrix}
a_1 & t & & \\
1 & a_2 & & \\
& 1 & \ddots & \\
& & \ddots & a_{n-1} & t\\
& & & 1 & a_n
\end{pmatrix}.
\end{equation}
换句话说，把 $M$ 送到 $(-a_1,\cdots,-a_n)$ 给出等同 $U\cong\mathbb{G}_m^n$。

现在用这个等同描述投影 $\mathrm{pr}_1$ 与 $\pi$。利用同构 $M$，可以把 $\star_1$ 的数据中的向量 $\{\mathbf v_i=e_{i+1}\}$、$\{\mathbf v_i=e_{i+1}\}$ 拉回，得到点 $\mathrm{pr}_1(M)=(\Lambda,\Lambda_*(\star_0),\mathbf v_i,\Lambda^*(\star_0),\mathbf v_i)\in\mathrm{Bun}^0_{n,1,2}$ 的相应向量：
\[
\mathbf v_i=a_{i+1}^{-1}e_{i+1},\quad i=0,\cdots,n-1;\qquad
\mathbf v_i=e_i-a_ie_{i-1},\quad i=1,\cdots,n.
\]
这意味着 (3.3) 中的 $M$ 对应的 $\mathrm{pr}_1(M)$ 的坐标为
\[
\mathrm{pr}_1(M)=j_0(a_1^{-1},\cdots,a_n^{-1},-a_1,\cdots,-a_n).
\]
另一方面，点 $\pi(M)\in\PP^1\setminus\{0,\infty\}$ 是使 $\det(M)=0$ 的 $t$ 的值，即
\[
\pi(M)=(-1)^na_1\cdots a_n.
\]
总之，$\mathrm{GR}^\circ_{\omega_1}$ 上的映射 $\mathrm{pr}_1$ 与 $\pi$ 可等同于如下映射
\[
\begin{tikzcd}[column sep=huge, row sep=large]
& \mathbb{G}_m^n \arrow[dl, "{(-\operatorname{inv},\iota)}"'] \arrow[dr, "\operatorname{mult}"] & \\
\mathbb{G}_m^n\times\mathbb{G}_a^n & & \mathbb{G}_m=\PP^1\setminus\{0,\infty\}
\end{tikzcd}
\]
其中 $\mathrm{inv}:\mathbb{G}_m^n\to\mathbb{G}_m^n$ 是逐坐标求逆，$\iota:\mathbb{G}_m^n\hookrightarrow\mathbb{G}_a^n$ 是自然包含。

于是得到

\begin{proposition}
与 Hecke 特征层 $A_{\varphi,\chi}$ 及对偶群 $GL_{n,\mathbb{Q}_\ell}=\widehat{GL}_n$ 的标准表示相伴的 Kloosterman 层 $\Kl^{\mathrm{Std}}_{GL_n}(\varphi,\chi)$ 形如
\[
\Kl^{\mathrm{Std}}_{GL_n}(\varphi,\chi)\cong\mathrm{mult}_!\big((-\mathrm{inv})^*\mathrm{Kum}_\chi\otimes\varphi^*\mathrm{AS}_\psi\big)[n-1].
\]
这里把 $\varphi:\mathbb{G}_a^n\to\mathbb{G}_a$ 到 $\mathbb{G}_m^n$ 的限制仍记为 $\varphi$。
\end{proposition}

\begin{remark}
当 $\varphi=\mathrm{add}:\mathbb{G}_a^n\to\mathbb{G}_a$ 是坐标求和，且把 $\chi$ 写成 $n$ 个乘性特征 $(\chi_1,\cdots,\chi_n)$ 时，Kloosterman 层 $\Kl^{\mathrm{Std}}_{GL_n}(\varphi,\chi)$ 与 Katz 在 [25, §4.1] 中定义的 Kloosterman 层 $\mathrm{Kl}(\psi;\chi_1,\cdots,\chi_n;1,\cdots,1)$ 相同，后者是 Deligne 的 Kloosterman 层 [10, §7] 的推广。
\end{remark}
